% Seminarhandout zur aktuellen main.tex.
% Mit XeLaTeX und BibTeX kompilieren: latexmk -xelatex handout.tex
\documentclass[11pt,a4paper,parskip=half]{scrartcl}
\usepackage{leipzig-handout}
% Angaben passend zur aktuellen Praesentation in TitlePageInfo.tex.
\newcommand{\HandoutTitel}{How to Vortrag und Feedback}
\newcommand{\HandoutUntertitel}{Fachliche Präsentationen richtig halten}
\newcommand{\HandoutKurztitel}{Vorträge und Feedback}
\newcommand{\HandoutAutor}{Jonas Bley}
\newcommand{\HandoutInstitut}{Professur für Physikdidaktik}
\newcommand{\HandoutVeranstaltung}{Fachdidaktik 1 Physik -- Seminar}
\newcommand{\HandoutDatum}{15.10.2026, WS 26/27}

\usepackage{csquotes}
\usepackage[backend=bibtex,style=authoryear,maxcitenames=2,
  maxbibnames=3,minbibnames=1,giveninits=true,uniquename=init,doi=true,url=false,isbn=false]{biblatex}
\addbibresource{references.bib}
\setlength{\bibitemsep}{3pt}
\setlength{\bibhang}{1em}
\renewcommand*{\bibfont}{\footnotesize}
\setlength{\emergencystretch}{1em}
\hypersetup{pdftitle={How to Vortrag und Feedback -- Handout},pdfauthor={Jonas Bley}}

% Gleichmaessige Schreibzeilen mit 8,5 mm Abstand.
\newcommand{\schreiblinien}[1]{%
  \par\begingroup\offinterlineskip\setlength{\parskip}{0pt}%
  \color{aquamarin!65!basalt}%
  \foreach \n in {1,...,#1}{\vbox to 8.5mm{\vfill\hrule width\linewidth height .4pt}}%
  \endgroup\par}

% Optional: \Randmusterfalse fuer eine besonders sparsame Druckfassung.
\begin{document}
\handouttitel

\section{Strategie und Vortragsstruktur}
Stimmt \textbf{Zielgruppe, Zielsetzung und Umsetzung} im jeweiligen Kontext aufeinander ab:
\begin{itemize}
  \item \textbf{Zielgruppe:} Was wissen die Zuhörenden bereits? Was erwarten sie?
  \item \textbf{Ziel:} Was sollen sie anschließend verstehen oder erklären können?
  \item \textbf{Umsetzung:} Welche fachliche Tiefe, Literatur, Sprache und
        Hilfsmittel passen zu Ziel, verfügbarer Zeit und Raum?
\end{itemize}

\begin{tabularx}{\linewidth}{@{}P{.16\linewidth}Y@{}}
  \toprule
  \textbf{Phase} & \textbf{Aufgabe} \\
  \midrule
  Einstieg & Interesse wecken, Vorwissen aktivieren, Ziel und Ablauf klären. \\
  Kernteil & Wenige klare Kernbotschaften mit Beispielen und Bildern erklären;
             Tiefe an die Zielgruppe anpassen und dem roten Faden folgen. \\
  Schluss & Kernaussagen aufgreifen, offene Fragen benennen, Fazit und Ausblick geben. \\
  \bottomrule
\end{tabularx}

\section{Folien, Visualisierungen und Hilfsmittel}
\begin{merksatz}[Gestaltungsziel]
Klare Aussagen und passende Abbildungen unterstützen eure Erklärung.
\end{merksatz}

\begin{description}[font=\normalfont\bfseries,leftmargin=0pt]
  \item[Kohärenz:] Unnötige Details und dekorative Bilder weglassen.
  \item[Signalisierung:] Den entscheidenden Zusammenhang hervorheben.
  \item[Räumliche Nähe:] Bild und zugehörige Beschriftung zusammen zeigen.
  \item[Dosierung:] Komplexe Erklärungen in nachvollziehbare Schritte teilen.
\end{description}
\quelle{\textcite{noetel2022,fiorella2021}.}

\textbf{Lesbarkeit:} Text kürzen, Kontraste und Freiraum nutzen und die Folie
aus der letzten Reihe prüfen.

\textbf{Hilfsmittel passend auswählen:} Folien, Handouts, Flipcharts,
Whiteboards und Overheadprojektoren; Modelle, Anschauungsmaterial und Experimente;
Videos, Animationen und Simulationen. Quizze, Umfragen und Abstimmungen ermöglichen
Beteiligung. Leitfrage: \emph{Was können die Zuhörenden damit besser verstehen
oder selbst erarbeiten?}

\newpage
\begingroup
\setlength{\parskip}{3pt plus .5pt minus .5pt}
\setlist{topsep=2pt,itemsep=1pt,parsep=0pt}
\section{Mündlich präsentieren}
\textbf{Du bist die Präsentation, nicht die Folien.}
Natürliche und effektive Kommunikation bedeutet:

\noindent\begin{minipage}[t]{.48\linewidth}
\textbf{Verbal}
\begin{itemize}
  \item Inhalte klar und verständlich vermitteln.
  \item Die Aufmerksamkeit der Zuhörenden halten.
  \item Botschaften betonen und verstärken.
\end{itemize}
\end{minipage}\hfill
\begin{minipage}[t]{.48\linewidth}
\textbf{Nonverbal und stimmlich}
\begin{itemize}
  \item Körpersprache, Gestik, Mimik und Blickkontakt bewusst einsetzen.
  \item Stimme, Lautstärke, Sprechtempo und Pausen variieren.
  \item Raum und Hilfsmittel effektiv nutzen.
\end{itemize}
\end{minipage}
\quelle{Nach \textcite{yates2002}.}

\section{Feedback im Lehrberuf}
\textbf{Feedback kann Lernen unterstützen. Sein Informationsgehalt ist entscheidend.}
\begin{itemize}
  \item \textbf{Lernerfolg:} 435 Studien zeigen einen mittleren Effekt von
        $d=0{,}48$ \parencite{10.3389/fpsyg.2019.03087}.
  \item \textbf{Unterrichtsqualität:} 23 Studien zu Rückmeldungen von
        Schüler*innen an Lehrkräfte zeigen eine kleine mittlere Verbesserung
        ($d=0{,}27$) \parencite{roehl2025}.
\end{itemize}
\textbf{Feedback für die Unterrichtsentwicklung nutzen:}
\begin{enumerate}
  \item \textbf{Interpretation unterstützen:} Rückmeldungen gemeinsam einordnen;
        intrinsische Motivation und positive Ermutigung fördern die Auseinandersetzung.
  \item \textbf{Dialog mit der Klasse führen:} Zu einzelnen Aspekten nachfragen
        und unterschiedliche Wahrnehmungen klären.
  \item \textbf{Umsetzung ermöglichen:} Gelegenheit schaffen, Rückmeldungen
        aufzugreifen und Veränderungen zu erproben.
\end{enumerate}

\subsection*{Feedback geben}
Rückmeldungen sollten \textbf{informativ, auf die Situation zugeschnitten und
handlungsorientiert} sein \parencite{brandmo2025}.
\begin{enumerate}
  \item \textbf{Beobachten, nicht bewerten:} Eine konkrete Stelle beschreiben:
        \enquote{Auf Folie X hast du das Thema von Argument Y bis Argument Z aufgearbeitet.}
  \item \textbf{Wirkung auf die eigene Person beschreiben:}
        \enquote{Für mich war das schwer verständlich und schwierig nachzuvollziehen.}
  \item \textbf{Handlungsorientiert sein:} Eine umsetzbare Anregung geben:
        \enquote{Du könntest Argument Z mit einem Beispiel untermauern oder
        die Argumentation mit einem Flussdiagramm beschreiben.}
  \item \textbf{Negativität und Schwammigkeit vermeiden:} Konkret und respektvoll formulieren.
  \item \textbf{Für Rückfragen bereit sein:} Beobachtungen und Vorschläge erläutern können.
\end{enumerate}

\subsection*{Feedback annehmen}
\textbf{Feedback nutzen heißt verstehen, abwägen und eine Verbesserung erproben.}
\begin{enumerate}
  \item \textbf{Bedanken:} Feedback als Geschenk wahrnehmen und die Rückmeldung wertschätzen.
  \item \textbf{Aktiv zuhören:} Die Rückmeldung vollständig aufnehmen und in eigenen Worten wiedergeben.
  \item \textbf{Nachfragen:} Um konkrete Beispiele bitten oder die zugrunde liegenden Schwierigkeiten klären.
  \item \textbf{Abwägen:} Die Rückmeldung mit dem Lernziel und der eigenen Beobachtung vergleichen.
  \item \textbf{Umsetzen:} Mögliche nächste Schritte beschreiben und Rückmeldung dazu einholen.
\end{enumerate}
\endgroup

\newpage
% Seite 3 ist die Vorderseite des Arbeitsbogens beim beidseitigen Druck.
\section{Anwendung: einen Optikvortrag planen}
Plant zu zweit anhand des bereitgestellten \textbf{Optik-Handouts} einen Vortrag
für eure Mitstudierenden. Haltet euren Plan hier fest. Auf der Rückseite ist Platz
für eure Visualisierungsideen.

\textbf{Ziel und Struktur des Vortrags}
\schreiblinien{7}
\vspace{3mm}
\textbf{Kernbotschaften für den Unterricht}
\schreiblinien{7}
\vspace{3mm}
\textbf{Methoden und Materialien}
\schreiblinien{7}

\newpage
% Seite 4 liegt direkt auf der Rueckseite des Planungsbogens.
\section*{Visualisierungsideen}
Haltet hier eure Ideen zur Visualisierung der \textbf{Vortragsstruktur}
oder der \textbf{Kernbotschaften} fest.

\vspace{4mm}
\begingroup
\setlength{\fboxsep}{0pt}
\setlength{\fboxrule}{.4pt}
\noindent\fcolorbox{aquamarin!65!basalt}{white}{%
  \parbox[t][212mm][t]{\dimexpr\linewidth-2\fboxrule\relax}{\strut}}
\endgroup

\newpage
\section{Feedbackmethoden im Überblick}
Die Frage bestimmt, welche Information ihr erhaltet. Rückmeldungen zum
Verständnis und Rückmeldungen zum Erleben der Veranstaltung erfüllen
unterschiedliche Zwecke.

{\small
\begin{tabularx}{\linewidth}{@{}P{.24\linewidth}Y@{}}
\toprule
\textbf{Methode} & \textbf{Durchführung und mögliche Fragen} \\
\midrule
5-Finger-Feedback & Handumriss beschriften: Daumen -- Das möchte ich weiterverfolgen
oder das fand ich gut; Zeigefinger -- Das habe ich gelernt; Mittelfinger -- Das hat
mich gestört; Ringfinger -- So war die Atmosphäre; kleiner Finger -- Das kam zu kurz. \\
Blitzlicht & Ein oder zwei Fragen stellen, z.\,B. „Was war heute besonders wichtig?“
oder „Was brauche ich noch?“ Alle antworten kurz und konkret in zwei bis drei Sätzen.
Während der Runde keine Kommentare oder Diskussion. \\
Digitale Umfrage & Beispielsweise mit Mentimeter Rückmeldungen anonym sammeln
und gemeinsam sichtbar machen: „Was hat euch gefallen?“ oder „Was könnte ich verbessern?“ \\
Feedbackanker & Vor dem Vortrag Beobachtungsaspekte verteilen, z.\,B. Inhalt,
Struktur, Körpersprache oder Stimme. Jede Person achtet auf ihren zugeteilten Aspekt.
Karten mit Farbe und Bezeichnung erleichtern die Zuordnung. \\
Interaktive Mindmap & Rückmeldungen gemeinsam nach Kategorien sammeln,
z.\,B. Inhalt, Struktur, Visualisierung und Verständlichkeit. \\
One-Minute-Paper & Eine kurze schriftliche Rückmeldung geben lassen:
„Was war der Kernpunkt?“ und „Welche Frage bleibt offen?“ Alternativ auf einer Karte
festhalten, was gefallen hat und welche Veränderung gewünscht ist. \\
\bottomrule
\end{tabularx}
}

\newpage
\section{Gesprächsbasierte Feedbackmethoden}
Diese beiden Abläufe geben den Vortragenden Raum, Rückfragen zu stellen,
die eigene Sicht einzubringen und auf Rückmeldungen einzugehen.

\subsection{Feedbackleiter mit Antwortphase}
\begin{enumerate}[itemsep=2pt]
  \item \textbf{Klären:} Das Publikum stellt Verständnisfragen;
        die vortragende Person beantwortet sie.
  \item \textbf{Würdigen:} Stärken mit konkreten Beispielen benennen.
  \item \textbf{Bedenken äußern:} Schwierigkeiten nachvollziehbar beschreiben.
  \item \textbf{Vorschlagen:} Ideen zur Verbesserung anbieten.
\end{enumerate}
\textbf{Ergänzte Antwortphase:} Die vortragende Person fasst die wichtigsten
Rückmeldungen in eigenen Worten zusammen, fragt bei Unklarheiten nach und
erläutert, welche Anregungen sie aufgreifen möchte und warum.
\quelle{\textcite{projectzeroLadder}}

\subsection{Feedbackgespräch nach dem R2C2-Modell}
\begin{enumerate}[itemsep=2pt]
  \item \textbf{Relationship -- Gesprächsbasis schaffen:} Gemeinsam einen
        Schwerpunkt wählen, etwa die Vortragsstruktur oder Verständlichkeit.
  \item \textbf{Reaction -- Reaktionen erkunden:} Eine Rückmeldung geben und
        die vortragende Person fragen, wie sie diese selbst sieht.
  \item \textbf{Content -- Inhalte klären:} Konkrete Situationen und
        unterschiedliche Wahrnehmungen gemeinsam besprechen.
  \item \textbf{Coaching -- nächste Schritte planen:} Gemeinsam vereinbaren,
        was erprobt werden kann. Dieser Schritt ist in der Seminarübung optional.
\end{enumerate}
\textbf{Reaktion der vortragenden Person:} Die eigene Sicht erläutern,
konkrete Beispiele erfragen und Vorschläge gemeinsam weiterentwickeln.
Rückmeldungen dürfen abgewogen werden; eine begründete andere Sicht ist möglich.
\quelle{\textcite{sargeant2017}}

\section{Literatur}
\printbibliography[heading=none]

\end{document}
